\appendix
\section{Proofs}\label{app:proofs}

\subsection{A price subgame without a pure equilibrium}\label{app:no-price-ne}

Consider the feasible history
\[
 x=\frac13,\qquad y=\frac23,\qquad a_1=a_2=0.
\]
Let $d=y-x=1/3$ and $q=p_1-p_2$. If $q<-d$, firm 1 is strictly preferred throughout the market; if $-d<q<d$, both firms serve positive shares and the indifferent consumer is $\hat z=(1-q)/2$; if $q>d$, firm 2 is strictly preferred throughout the market. At $q=\pm d$, a positive-measure backyard segment is indifferent, so the argument must allow arbitrary tie allocation.

Suppose first that $-d<q<d$. Both prices must be strictly positive in equilibrium: a positive-demand firm charging zero can raise its price slightly and earn positive profit while remaining in the same branch. The first-order conditions on this smooth branch are
\[
 p_1=\frac{1+p_2}{2},
 \qquad
 p_2=\frac{1+p_1}{2},
\]
whose unique solution is $(p_1,p_2)=(1,1)$. Each firm then earns normalized revenue $1/2$. But firm 1 can choose $p_1'=7/12$. Since
\[
 \frac7{12}<1-\frac13=\frac23,
\]
firm 1 becomes strictly preferred even at $z=1$, captures the entire market, and earns $7/12>1/2$. Hence there is no equilibrium with $|q|<d$.

If $q<-d$, firm 2 has zero demand and necessarily $p_2>p_1+d>0$. It can reduce its price to a strictly positive number below $p_1+d$ and obtain positive demand and positive profit. Hence this region contains no equilibrium. The case $q>d$ is symmetric.

It remains to consider $q=-d$. Write $p_2=P$ and $p_1=P-d$. The right backyard segment of length $1/3$ is indifferent. Let $\tau\in[0,1]$ denote the fraction of that segment assigned to firm 1. If $\tau>0$, firm 2 can undercut by an arbitrarily small $\varepsilon>0$, move to the interior branch, and obtain a limiting gain $P\tau/3>0$. If $\tau=0$ and $P>d$, firm 1 can undercut slightly, capture the full market, and obtain a limiting gain $(P-d)/3>0$. Finally, if $\tau=0$ and $P=d$, then $p_1=0$ and firm 1 can instead raise its price to a sufficiently small $\varepsilon>0$, obtaining
\[
 \varepsilon\left(\frac23-\frac{\varepsilon}{2}\right)>0.
\]
Thus no tie allocation sustains $q=-d$. The case $q=d$ is symmetric. The price subgame therefore has no pure Nash equilibrium. Since this is a proper subgame after a feasible history, a pure-strategy SPNE cannot exist under the standard modern full-history definition. \qed

\subsection{Proof of Lemma~\ref{lem:price}}\label{app:price-proof}

At maximal locations, if both firms serve positive shares, firm 1's market share is
\[
 s_1=\frac{1+\delta-p_1+p_2}{2},
\]
and $s_2=1-s_1$. On the interior branch, each firm's revenue is strictly concave in its own price. The first-order conditions yield \eqref{eq:regular-prices}. At those prices, $s_i=p_i/(2)$, so normalized revenue is $p_i^2/2$.

A deviation by firm 1 that captures the whole market must satisfy
\[
 p_1\leq p_2+\delta-1.
\]
Evaluated against the interior candidate $p_2=1-\delta/3$, the capture threshold is $2\delta/3$. Because prices are constrained to be nonnegative, a feasible nonnegative full-market capture price for firm 1 exists on the regular branch only when $\delta\geq0$. For $0\leq\delta\leq3$, the highest such capture price is $2\delta/3$, and the difference between firm 1's interior-candidate revenue and the feasible capture revenue is
\[
 \frac12\left(1+\frac{\delta}{3}\right)^2-\frac{2\delta}{3}
 =\frac{(\delta-3)^2}{18}\geq0.
\]
For $-3\leq\delta<0$, the capture threshold is negative, so the nonnegative-price strategy set contains no firm-1 full-market capture deviation against the interior candidate. Symmetrically, firm 2's capture threshold is $-2\delta/3$: a feasible nonnegative full-market capture price exists only when $\delta\leq0$, and for $-3\leq\delta\leq0$ the corresponding comparison is $(\delta+3)^2/18\geq0$. For $0<\delta\leq3$, firm 2's full-market capture set is empty under the nonnegative-price constraint.

Together with strict concavity on the shared-market branch and the fact that a move into a zero-demand region yields zero revenue, these feasible capture comparisons rule out profitable deviations from the interior candidate for $-3<\delta<3$. At $\delta=0$ both capture thresholds equal zero and the square comparisons are consistent with the shared-market optimum. At $\delta=3$ (respectively $\delta=-3$), the regular candidate coincides with the firm-1 (respectively firm-2) exclusion boundary, so the limiting branch formulas agree. Thus the interior candidate is globally stable on its valid regular branch without treating a negative capture threshold as a feasible price deviation.

Now suppose $\delta>3$. Against $p_2=0$, firm 1 can capture the whole market at any $p_1\leq\delta-1$, so its best full-market price is $p_1=\delta-1$. For $p_1\geq\delta-1$, its interior-branch revenue is
\[
 r_1(p_1)=\frac{p_1(1+\delta-p_1)}{2}.
\]
The unconstrained maximizer $(1+\delta)/2$ lies strictly below $\delta-1$ when $\delta>3$, so revenue decreases as firm 1 moves above the capture boundary. Thus $p_1=\delta-1$ is a best response to $p_2=0$. Firm 2 earns zero. Any positive $p_2$ paired with a capture price for firm 1 cannot form an equilibrium because firm 2 can undercut into positive demand or firm 1 can raise its capture price. Hence the equilibrium pair is $(\delta-1,0)$. The case $\delta<-3$ is symmetric. At $\delta=\pm3$ the formulas coincide. \qed

\subsection{Proof of Proposition~\ref{prop:correction}}\label{app:correction-proof}

Fix the rival's quality at
\[
 A=\frac1{3\lambda}
\]
and restrict attention to $\lambda\in(1/9,2/9)$. Since $A<3$ throughout this interval, a nonnegative deviation by firm 1 never enters the low-quality exclusion region $a_1-A\leq-3$. On the regular branch, $0\leq a_1\leq A+3$, normalized profit is
\[
 f(a_1)=\frac12\left(1+\frac{a_1-A}{3}\right)^2-\frac{\lambda a_1^2}{2}.
\]
Because $f''(a_1)=1/9-\lambda<0$ and $f'(A)=0$, the unique maximizer of the regular branch is $a_1=A$.

On the high-quality exclusion branch, $a_1\geq A+3$, normalized profit is
\[
 g(a_1)=a_1-A-1-\frac{\lambda a_1^2}{2}.
\]
This branch is strictly concave and is maximized at $a_1=1/\lambda$. Moreover,
\[
 \frac1\lambda-A=\frac{2}{3\lambda}>3
\]
for $\lambda<2/9$, so the maximizer lies in the exclusion branch. The regular and exclusion maxima are
\[
 f(A)=\frac12-\frac1{18\lambda},
 \qquad
 g(1/\lambda)=\frac1{6\lambda}-1.
\]
Their difference is
\[
 g(1/\lambda)-f(A)=\frac{4-27\lambda}{18\lambda}.
\]
Hence $A$ fails to be a best response when $\lambda<4/27$, ties with the exclusionary action at $\lambda=4/27$, and is the global best response for $4/27<\lambda<2/9$. Symmetry completes the proof. \qed

\subsection{Proof of Proposition~\ref{prop:asymmetric}}\label{app:asymmetric-proof}

First fix the rival's quality at zero. For $0\leq a\leq3$, normalized profit is
\[
 f_0(a)=\frac12\left(1+\frac a3\right)^2-\frac{\lambda a^2}{2},
\]
with derivative
\[
 f_0'(a)=\frac13+a\left(\frac19-\lambda\right).
\]
At $a=3$,
\[
 f_0'(3)=\frac{2-9\lambda}{3}>0
\]
for $\lambda<2/9$, so the regular-branch payoff is increasing through the regime boundary. On the exclusion branch $a\geq3$,
\[
 g_0(a)=a-1-\frac{\lambda a^2}{2},
\]
which is strictly concave and maximized at $a=1/\lambda>3$. Hence
\[
 BR(0)=\frac1\lambda.
\]

Now fix the rival's quality at $B=1/\lambda$. On the region $0\leq a\leq B-3$, the firm is excluded and earns $-\lambda a^2/2$, uniquely maximized at $a=0$. On the regular region $a\in[B-3,B+3]$, the derivative vanishes only at
\[
 a^{c}=\frac{3\lambda-1}{\lambda(9\lambda-1)}<0,
\]
so the regular-branch payoff is decreasing throughout its feasible positive domain. Its value at the lower boundary is
\[
 -\frac{\lambda(B-3)^2}{2}<0.
\]
Finally, on the reverse-exclusion branch $a\geq B+3$, profit is
\[
 a-B-1-\frac{\lambda a^2}{2}.
\]
Its unconstrained maximizer is $1/\lambda=B$, below the branch, so the payoff is decreasing on the entire reverse-exclusion region and is negative at $B+3$. Thus zero quality is the unique best response to $B$:
\[
 BR(1/\lambda)=0.
\]
The profile $(1/\lambda,0)$ is therefore a Nash equilibrium, and symmetry gives $(0,1/\lambda)$. \qed
